\subsection{Componentes y configuración de CP2}
La copia de código preparada en \texttt{checkpoint2/Checkpoint2\_LDR/}
usa ESP32, WiFi y PubSubClient; lee GPIO35, imprime a 115200 baudios y
publica el valor crudo como texto en \texttt{lab/ldr}. El bucle contiene
una espera de 1000~ms: no representa una medición de período exacto,
pues las operaciones de red también consumen tiempo. La configuración
local se mantiene separada del código versionado.

\pendiente{Confirmar modelo de placa, versiones efectivamente usadas,
resistor, orientación del divisor y alimentación. Describir montaje real
y pasos ejecutados, con fotografías, sin completar valores por suposición.}
\pendiente{CP1: completar guía breve de instalación y carga de prueba
con los pasos realmente seguidos, incidencias y sus capturas.}

\evidencia{B1}{Insertar diagrama de bloques del montaje real de CP2.
Propuesta a validar: LDR/divisor, ADC ESP32, broker y receptor en otro equipo.}
{Arquitectura de adquisición y distribución MQTT de CP2 (diagrama pendiente).}
\evidencia{F1}{Insertar flujo correspondiente al código final confirmado:
conexión WiFi, lectura, conexión MQTT, servicio del cliente y publicación;
incluir la rama de reconexión fallida.}
{Flujo del programa de telemetría de CP2 (diagrama pendiente).}

\subsection{Resto de experiencias}
\pendiente{Integrar materiales y procedimiento de CP3/CP4 desde el registro
del equipo y su sesión responsable. No reconstruir una ejecución a partir
de los resultados esperados de la guía.}
